\documentclass[11pt,a4paper]{article}

\usepackage[T1]{fontenc}
\usepackage[utf8]{inputenc}
\usepackage[spanish,es-nodecimaldot]{babel}
\usepackage{lmodern}
\usepackage{microtype}
\usepackage[a4paper,margin=2.4cm]{geometry}
\usepackage{amsmath}
\usepackage{amssymb}
\usepackage{mathtools}
\usepackage{booktabs}
\usepackage{longtable}
\usepackage{array}
\usepackage{tabularx}
\usepackage{hyperref}
\usepackage{enumitem}
\usepackage{xcolor}

\hypersetup{
  colorlinks=true,
  linkcolor=blue!50!black,
  urlcolor=blue!50!black,
  pdftitle={Estado actual del proyecto Halo2 wrapper},
  pdfauthor={OpenAI Codex}
}

\setlist[itemize]{leftmargin=1.4em}
\setlist[enumerate]{leftmargin=1.6em}
\setlength{\parskip}{0.5em}
\setlength{\parindent}{0pt}
\renewcommand{\arraystretch}{1.15}

\title{Documentaci\'on extensiva del estado actual del proyecto\\
\large Halo2 wrapper para verificaci\'on Groth16 BN254 en un outer circuit Halo2/Midnight}
\author{Documento t\'ecnico interno}
\date{Estado relevado del repositorio al 28 de mayo de 2026}

\begin{document}

\maketitle
\tableofcontents
\newpage

\section{Prop\'osito, alcance y advertencias}

Este documento describe el estado real del repositorio \texttt{halo2-groth16-bn254-to-bls12-381}
tal como est\'a implementado hoy. El objetivo no es presentar una visi\'on idealizada del proyecto,
sino fijar con precisi\'on:

\begin{itemize}
  \item qu\'e piezas criptogr\'aficas, circuitales y de backend existen hoy;
  \item qu\'e optimizaciones locales ya fueron retenidas, intentadas o descartadas;
  \item cu\'al es el cuello de botella operativo actual;
  \item cu\'al es el camino m\'as directo para destrabar la lane de ejecuci\'on real.
\end{itemize}

El repositorio est\'a en \textbf{Stage 1 / Week 5+} y ya no es una simple base de workspace:
incluye una lane directa real de \texttt{setup -> prove -> verify} sobre el circuito can\'onico
\texttt{OuterWrapperCircuit}. Sin embargo, \textbf{no} debe leerse como un wrapper verifier
productivo, general o endurecido para producci\'on.

La slice actual es deliberadamente estrecha:

\begin{itemize}
  \item el objetivo principal es envolver verificaci\'on Groth16 BN254 en una estructura outer Halo2/Midnight;
  \item la aritm\'etica BN254 est\'a implementada s\'olo hasta la frontera necesaria para el pairing core y la slice Groth16 actual;
  \item la lane directa existe, pero sigue bloqueada por presi\'on de memoria en la etapa \texttt{prove-finalize}.
\end{itemize}

\subsection{Fuente de verdad}

Ante conflictos entre documentos, la fuente de verdad debe considerarse en este orden:

\begin{enumerate}
  \item el c\'odigo vigente del repositorio;
  \item los ADRs aceptados bajo \texttt{docs/decisions/};
  \item los documentos operativos y de profiling bajo \texttt{docs/};
  \item referencias heredadas en \texttt{README.md} y \texttt{AGENTS.md} que puedan haber quedado desactualizadas.
\end{enumerate}

Este criterio es especialmente importante porque el repositorio hoy contiene algunas derivas
documentales: hay nombres viejos de archivos, rutas renombradas y una referencia faltante
a un plan de estrategia del prover que ya no est\'a presente en el \'arbol de archivos.

\subsection{Tabla de alcance actual}

\begin{longtable}{>{\raggedright\arraybackslash}p{0.22\textwidth} >{\raggedright\arraybackslash}p{0.72\textwidth}}
\toprule
\textbf{Categor\'ia} & \textbf{Contenido} \\
\midrule
\endhead

\textbf{Implementado hoy} &
Workspace multi-crate; parsing real de artefactos \texttt{snarkjs}; capa BN254 en circuito para
\texttt{Fp}, \texttt{Fp2}, \texttt{Fp6}, \texttt{Fp12}, G1, G2 af\'in, G2 Jacobian estrecho,
Miller path, final exponentiation, pairing-product check estrecho, slice Groth16 BN254,
\texttt{OuterWrapperCircuit}, backend directo Halo2/Midnight, fixtures \texttt{circom\_multiplier2}
y \texttt{semaphore}, profiling de layout y contratos de setup/prove/verify. \\

\textbf{Parcialmente implementado / experimental} &
Split \texttt{prove-trace} / \texttt{prove-finalize}; artefacto de setup enriquecido; patch local
de \texttt{midnight-proofs}; chunking y materializaci\'on perezosa parcial en \texttt{compute\_h\_poly(...)}; lane can\'onica R1CS como backend alterno futuro; exploraci\'on de kernels torus/compressed-region. \\

\textbf{Expl\'icitamente fuera de alcance actual} &
Verificador Groth16 general; APIs p\'ublicas amplias de pairing o multi-pairing; subgroup checks de G2;
scalar multiplication amplia en G2; wrapper verifier productivo; hardening de seguridad de producci\'on;
CI r\'apido y estable para toda la lane pesada; optimizaci\'on global del prover m\'as all\'a de la lane estrecha hoy existente. \\

\bottomrule
\end{longtable}

\section{Arquitectura del workspace y flujo de datos}

\subsection{Responsabilidades por crate}

\textbf{\texttt{wrapper-core}} concentra conceptos de dominio estables y evita, en lo posible,
dependencias Halo2. All\'i viven \texttt{WrapperJob}, \texttt{WrapperExecutionPackage},
\texttt{WrapperExecutionResult}, modelado de statement, packaging y contratos de artefactos.

\textbf{\texttt{wrapper-circuits}} contiene la sem\'antica circuital y la capa BN254 activa:
tipos asignados, gadgets sobre \texttt{midnight-circuits}, la slice Groth16 BN254 y el
\texttt{OuterWrapperCircuit}. Tambi\'en aloja la l\'inea R1CS can\'onica, pero \'esta no es hoy
el camino cr\'itico de entrega.

\textbf{\texttt{wrapper-backends}} encapsula parsing, adaptaci\'on, normalizaci\'on de artefactos y
la lane directa de outer proofs. Su responsabilidad es transformar artefactos externos
en una forma apta para la capa circuital sin redefinir la sem\'antica del circuito.

\textbf{\texttt{wrapper-cli}} expone comandos honestos de inspecci\'on, profiling, setup, prove,
verify y diagn\'ostico.

\textbf{\texttt{wrapper-tests}} agrupa fixtures, helpers de integraci\'on y benchmarks.

\subsection{Flujo de datos real actual}

La lane principal implementada hoy, y la que conviene priorizar operativamente, es:

\begin{quote}
\texttt{snarkjs artifacts}\\
\(\downarrow\)\\
\texttt{Groth16Bn254ArtifactBundle}\\
\(\downarrow\)\\
\texttt{WrapperJob}\\
\(\downarrow\)\\
\texttt{WrapperExecutionPackage}\\
\(\downarrow\)\\
\texttt{OuterWrapperCircuit}\\
\(\downarrow\)\\
\texttt{MidnightDirectOuterBackend}\\
\(\downarrow\)\\
\texttt{setup/prove/verify}
\end{quote}

En forma m\'as concreta:

\begin{enumerate}
  \item \texttt{wrapper-backends/src/snarkjs.rs} parsea \texttt{proof.json}, \texttt{public.json}
  y \texttt{verification\_key.json} en formato \texttt{snarkjs bn128}.
  \item \texttt{wrapper-backends/src/groth16.rs} normaliza eso en
  \texttt{Groth16Bn254ArtifactBundle}.
  \item Ese bundle se transforma en \texttt{WrapperJob} y luego en
  \texttt{WrapperExecutionPackage}.
  \item \texttt{wrapper-circuits/src/outer/} adapta el package a la entrada can\'onica del outer circuit.
  \item \texttt{wrapper-backends/src/outer/direct/proving.rs} ejecuta la lane real de setup/prove/verify.
\end{enumerate}

\subsection{Separaci\'on entre sem\'antica y backend}

La decisi\'on arquitect\'onica relevante es que el circuito can\'onico vive en
\texttt{wrapper-circuits}, mientras que \texttt{wrapper-backends} expresa la lane concreta de
materializaci\'on de artefactos. Esto evita que parsing, serializaci\'on o detalles del prover
reemplacen la definici\'on sem\'antica del wrapper.

\section{Fundamentos matem\'aticos}

\subsection{Campo base y torre de extensi\'on}

El proyecto trabaja sobre la familia BN254. La implementaci\'on usa la torre compatible con
arkworks/Midnight:

\[
\mathbb{F}_{p^2} = \mathbb{F}_p[u]/(u^2 + 1)
\]

\[
\mathbb{F}_{p^6} = \mathbb{F}_{p^2}[v]/(v^3 - (9 + u))
\]

\[
\mathbb{F}_{p^{12}} = \mathbb{F}_{p^6}[w]/(w^2 - v)
\]

Las convenciones de representaci\'on activas en el c\'odigo son:

\begin{itemize}
  \item \(\texttt{Fp2} = c_0 + c_1 u\), con orden de coordenadas \((c_0, c_1)\);
  \item \(\texttt{Fp6} = c_0 + c_1 v + c_2 v^2\), con orden \((c_0, c_1, c_2)\);
  \item \(\texttt{Fp12} = c_0 + c_1 w\), con orden \((c_0, c_1)\).
\end{itemize}

\subsection{Curvas, twist y representaciones}

La slice actual necesita:

\begin{itemize}
  \item G1 sobre \(\mathbb{F}_p\);
  \item G2 sobre el twist cuadr\'atico sobre \(\mathbb{F}_{p^2}\);
  \item una representaci\'on af\'in estrecha en G2;
  \item una representaci\'on Jacobian estrecha para ciertos pasos;
  \item una representaci\'on homog\'enea preparada para Miller traversal.
\end{itemize}

En G2 se implementan validaciones de curva sobre el twist, pero no subgroup checks amplios.

\subsection{Funciones de l\'inea y embedding disperso}

Durante el Miller path, los coeficientes de l\'inea se exponen en la forma dispersa:

\[
(\ell_0,\ell_w,\ell_{vw})
\]

Para un punto af\'in \(P = (x_P, y_P)\) de G1, la evaluaci\'on correspondiente es:

\[
\ell_0 y_P + \ell_w x_P w + \ell_{vw} v w
\]

La documentaci\'on y el c\'odigo convergen en que este embedding disperso mapea naturalmente a
slots de \(\mathbb{F}_{p^{12}}\) compatibles con la ruta tipo \texttt{mul\_by\_034}.
El borde p\'ublico activo para consumir esta estructura es
\texttt{AssignedMillerAccumulator::mul\_by\_line(...)}.

\subsection{Miller loop}

La lane de pairing implementada es una optimal-ate Miller traversal con schedule preparado y
determinista. El repositorio no expone hoy una API p\'ublica amplia de pairings; expone un
pairing-product check estrecho y verificador-shaped.

\subsection{Final exponentiation}

La final exponentiation sigue la descomposici\'on cl\'asica en:

\begin{itemize}
  \item \textbf{easy part}: componentes de bajo costo relativo que usan conjugaci\'on, inversi\'on y Frobenius;
  \item \textbf{hard part}: principal hotspot de filas, donde domina la rutina \texttt{exp\_by\_neg\_x(...)}.
\end{itemize}

La hard part es el foco principal de optimizaci\'on local del pairing core, porque impacta:

\begin{itemize}
  \item \texttt{bn254\_final\_exponentiation\_hard\_part};
  \item \texttt{bn254\_final\_exponentiation};
  \item \texttt{bn254\_pairing\_check\_sample};
  \item \texttt{bn254\_pairing\_check\_groth16\_style};
  \item el verificador Groth16 end-to-end.
\end{itemize}

\section{Estado implementado del stack criptogr\'afico}

\subsection{Capa de campo y extensiones}

La capa BN254 activa en circuito cubre:

\begin{itemize}
  \item \texttt{AssignedFp}: valores en \(\mathbb{F}_p\) sobre chips de foreign field;
  \item \texttt{AssignedFp2}: pares de \texttt{AssignedFp};
  \item \texttt{AssignedFp6}: ternas de \texttt{AssignedFp2};
  \item \texttt{AssignedFp12}: pares de \texttt{AssignedFp6}.
\end{itemize}

Las operaciones activas son, seg\'un el nivel:

\begin{itemize}
  \item \texttt{zero}, \texttt{one}, \texttt{add}, \texttt{sub}, \texttt{neg}, igualdad;
  \item \texttt{mul} y \texttt{square};
  \item helpers internos para no-residue, conjugaci\'on, Frobenius y kernels usados por pairing.
\end{itemize}

\subsection{Capa de curva}

En G1 existe soporte m\'inimo real para:

\begin{itemize}
  \item asignaci\'on;
  \item construcci\'on desde coordenadas con chequeo \emph{on-curve};
  \item suma;
  \item cobertura de tests con referencia arkworks.
\end{itemize}

En G2 existe hoy:

\begin{itemize}
  \item \texttt{AssignedG2Affine} con asignaci\'on, negaci\'on, igualdad y \emph{on-curve};
  \item \texttt{AssignedG2Projective} en Jacobian estrecho con \texttt{from\_affine}, \texttt{neg},
  \texttt{double} y \texttt{add} incompleta;
  \item \texttt{AssignedG2MillerPoint} con \texttt{double\_with\_line} y
  \texttt{mixed\_add\_with\_line}.
\end{itemize}

\subsection{Pairing core}

La slice Week 4 efectivamente aterrizada incluye:

\begin{itemize}
  \item traversal real de Miller;
  \item acumulaci\'on en \(\mathbb{F}_{p^{12}}\);
  \item final exponentiation;
  \item pairing-product check que multiplica primero los outputs de Miller y aplica
  exactamente \emph{una} final exponentiation compartida.
\end{itemize}

Esta decisi\'on es importante: el repositorio no realiza final exponentiation por t\'ermino, sino una
reducci\'on verificador-shaped acorde al producto total.

\subsection{Slice Groth16 BN254}

El m\'odulo \texttt{crates/wrapper-circuits/src/groth16.rs} implementa una slice estrecha de
verificaci\'on Groth16 BN254 con las siguientes restricciones deliberadas:

\begin{itemize}
  \item el \texttt{proof} y la \texttt{verification key} ya llegan normalizados;
  \item los public inputs s\'olo se usan como escalares del acumulador de IC;
  \item la verificaci\'on se reduce directamente a un pairing-product check;
  \item no se presenta un framework Groth16 general ni multi-backend.
\end{itemize}

\subsection{Outer circuit y backend directo}

El circuito outer can\'onico es \texttt{OuterWrapperCircuit}. Su sem\'antica es:

\begin{quote}
verificar la slice Groth16 BN254 soportada hoy y exponer el statement outer correspondiente.
\end{quote}

El backend directo concreto es \texttt{MidnightDirectOuterBackend}, con variantes de host lane.
La implementaci\'on activa permite:

\begin{itemize}
  \item setup real;
  \item prove real;
  \item verify real;
  \item split adicional \texttt{prove-trace} / \texttt{prove-finalize}.
\end{itemize}

\section{Verificador Groth16 BN254 actual}

\subsection{Ecuaci\'on de verificaci\'on usada}

La reducci\'on concreta implementada es:

\[
e(A,B) \cdot e(-vk_x,\gamma) \cdot e(-C,\delta) = e(\alpha,\beta)
\]

Aqu\'i:

\begin{itemize}
  \item \(A\) y \(C\) est\'an en G1;
  \item \(B\) est\'a en G2;
  \item \((\alpha,\beta,\gamma,\delta)\) provienen de la verification key;
  \item \(vk_x\) es la combinaci\'on lineal del vector IC con los public inputs.
\end{itemize}

\subsection{Acumulaci\'on de \texorpdfstring{$vk_x$}{vk\_x}}

Si la verification key trae puntos
\(\texttt{ic}[0], \texttt{ic}[1], \dots, \texttt{ic}[n]\)
y los public inputs son
\(x_1,\dots,x_n\),
entonces:

\[
vk_x = \texttt{ic}[0] + \sum_{i=1}^{n} x_i \cdot \texttt{ic}[i]
\]

En el repositorio actual esta acumulaci\'on es:

\begin{itemize}
  \item verifier-only;
  \item apoyada sobre la G1 chip existente;
  \item estrecha y no dise\~nada como una API p\'ublica de MSM general.
\end{itemize}

\subsection{Camino variable y camino preparado constante}

La shape actual del pairing check Groth16 est\'a optimizada de manera verificador-shaped:

\begin{itemize}
  \item el t\'ermino de prueba \(e(A,B)\) conserva camino variable sobre G2;
  \item el t\'ermino completamente fijo \(e(\alpha,\beta)\) se precomputa
  off-circuit como una constante de GT;
  \item los t\'erminos constantes restantes de verification key
  \(\gamma_{g2}\) y \(\delta_{g2}\)
  se preprocesan off-circuit en una forma preparada constante de Miller.
\end{itemize}

En otras palabras, el circuito evita pagar el costo completo de recorrer el
t\'ermino \((\alpha,\beta)\) dentro del Miller loop y tampoco recorre
\(\gamma_{g2}\) ni \(\delta_{g2}\) como si fueran witness variables. Esta es
una de las decisiones estructurales m\'as valiosas ya retenidas en la slice
Groth16 actual.

\section{Historial de optimizaciones: retenidas, intentadas y descartadas}

\subsection{Tabla consolidada}

\scriptsize
\begin{longtable}{>{\raggedright\arraybackslash}p{0.14\textwidth} >{\raggedright\arraybackslash}p{0.10\textwidth} >{\raggedright\arraybackslash}p{0.16\textwidth} >{\raggedright\arraybackslash}p{0.16\textwidth} >{\raggedright\arraybackslash}p{0.21\textwidth} >{\raggedright\arraybackslash}p{0.15\textwidth}}
\toprule
\textbf{Cambio} & \textbf{Estado} & \textbf{Motivaci\'on} & \textbf{Ubicaci\'on conceptual} & \textbf{M\'etrica antes/despu\'es} & \textbf{Conclusi\'on} \\
\midrule
\endhead

\texttt{mul\_by\_constant(...)} & Retenida & Reemplazar multiplies gen\'ericas por constantes chicas con una ruta m\'as barata del foreign-field chip & Helpers de torre BN254, especialmente \texttt{mul\_by\_nonresidue\_fp2} & Se documenta como una de las principales razones por las que el verificador total baj\'o de \texttt{k=22} a \texttt{k=21} & Es la primitive local m\'as confiable descubierta hasta ahora \\

\texttt{add\_constant(...)} en \texttt{AssignedG2Affine::assert\_on\_curve(...)} & Retenida & Evitar witness fijo para el coeficiente de twist y absorber el offset como constante del chip & G2 af\'in y preparaci\'on para Miller & \texttt{g2 on\_curve: 400 -> 378}; \texttt{g2 neg: 930 -> 886}; \texttt{g2 proj double: 2594 -> 2550}; \texttt{g2 mixed\_add\_with\_line: 3374 -> 3330} & Win local real en G2/prep, sin mover por s\'i sola el total del pairing core \\

Signed-window \texttt{exp\_by\_neg\_x(...)} & Retenida & Cambiar una cadena positiva por una signed chain que ahorra un multiply principal por llamada & Hard part de final exponentiation & \texttt{hard part: 574112 -> 561254}; \texttt{final exponentiation: 587420 -> 574562}; \texttt{pairing sample: 1682524 -> 1669666}; \texttt{Groth16-style: 1949238 -> 1936380} & Win retenida y can\'onica; la metadata vive en \texttt{final\_exp\_chain.rs} \\

Compressed cyclotomic squaring dentro de \texttt{exp\_by\_neg\_x(...)} & Retenida & Aprovechar que los bloques repetidos de squaring tienen una estructura mejor que el cuadrado ciclot\'omico gen\'erico & Repeated square blocks de la hard part & \texttt{hard part: 561254 -> 492083}; \texttt{final exponentiation: 574562 -> 505391}; \texttt{pairing sample: 1669666 -> 1600495}; \texttt{Groth16-style: 1936380 -> 1867209} & Es la mejor optimizaci\'on local retenida hoy en el pairing core \\

Precomputar \texttt{e(\alpha,\beta)} como constante de GT & Retenida & Sacar del multi-Miller loop el t\'ermino de verification key que es completamente fijo & Reducci\'on Groth16 verifier-shaped & \texttt{pairing check Groth16-style: 1866759 -> 1632579} & Win estructural fuerte: cuando el t\'ermino es totalmente constante, colapsarlo fuera del Miller loop gana bastante m\'as que dejarlo s\'olo como G2 preparado \\

\texttt{linear\_combination(...)} sobre \texttt{Fp2/Fp6/Fp12} hot path & Descartada & Hip\'otesis de que una forma af\'in compacta reducir\'ia filas & \texttt{AssignedFp2::mul\_by\_constant}, \texttt{AssignedFp6::mul\_by\_nonresidue\_fp2}, helpers \texttt{3t \textpm 2z} & \texttt{fp6 mul: 1252 -> 1318}; \texttt{fp12 cyclotomic square: 1622 -> 1886}; \texttt{final exponentiation: 587420 -> 678119}; \texttt{pairing check: 1682524 -> 1805233} & Regresi\'on medida; no debe relandearse salvo con shape materialmente distinta \\

Helpers \texttt{select}/\texttt{is\_equal*}/\texttt{is\_zero} para identidad GT & Descartada como optimizaci\'on & Limpiar control de flujo y empaquetar la comparaci\'on final de identidad & Check final del pairing-product & Row-neutral seg\'un \texttt{doctor} y \texttt{profile-layout --family blocks} & Puede servir para claridad, no para rendimiento medido \\

Prototipo torus-style aislado en \texttt{y7} & Descartada & Explorar \texttt{cyclotomic * unitary\_inverse(cyclotomic)} mediante coordenada de torus & Sitio aislado de la hard part & \texttt{hard part: 561254 -> 571604}; \texttt{final exponentiation: 574562 -> 584912}; \texttt{pairing sample: 1669666 -> 1680016}; \texttt{Groth16-style: 1936380 -> 1946730} & Una sustituci\'on aislada no amortiza compresi\'on/descompresi\'on \\

Rollout de \texttt{CyclotomicFp12MulChip} & Descartada & Reempaquetar multiplicaciones ciclot\'omicas repetidas como gadget expl\'icito & Sitios \texttt{y3}, \texttt{y9}, \texttt{y10}, \texttt{y11} & \texttt{hard part: 561254 -> 561344}; \texttt{final exponentiation: 574562 -> 574652}; \texttt{pairing sample: 1669666 -> 1669756}; \texttt{Groth16-style: 1936380 -> 1936470} & Reempaque sin kernel nuevo no alcanza para ganar filas \\

Refactors schedule-only del Miller accumulator & Poco prometedor / no retenido & Intentar ganar por ordenamiento o fusi\'on superficial & Secuenciaci\'on de square + line consumption & Sin reducci\'on material de filas & La siguiente mejora debe atacar costo aritm\'etico subyacente, no wrapper sequencing \\

\bottomrule
\end{longtable}
\normalsize

\subsection{Cadena retenida de \texttt{exp\_by\_neg\_x(...)}}

La cadena signed-window can\'onica actual parte de \(35\) y consume:

\begin{itemize}
  \item \texttt{\string<\string< 6, -35}
  \item \texttt{\string<\string< 9, +101}
  \item \texttt{\string<\string< 8, -83}
  \item \texttt{\string<\string< 9, +37}
  \item \texttt{\string<\string< 9, +105}
  \item \texttt{\string<\string< 11, +79}
  \item \texttt{\string<\string< 5, +17}
\end{itemize}

La metadata compartida entre host y circuito vive en
\texttt{crates/wrapper-circuits/src/bn254/final\_exp\_chain.rs}.

\subsection{Tabla de mediciones clave}

\begin{tabularx}{\textwidth}{>{\raggedright\arraybackslash}p{0.42\textwidth} >{\raggedright\arraybackslash}X}
\toprule
\textbf{Bloque / m\'etrica} & \textbf{Serie relevante documentada} \\
\midrule
\texttt{g2 on\_curve} & \texttt{400 -> 378} \\
\texttt{final exponentiation hard part} & \texttt{574112 -> 561254 -> 492083} \\
\texttt{final exponentiation} & \texttt{587420 -> 574562 -> 505391} \\
\texttt{pairing check sample} & \texttt{1682524 -> 1669666 -> 1600495} \\
\texttt{pairing check Groth16-style} & \texttt{1949238 -> 1936380 -> 1867209 -> 1866759 -> 1632579} \\
\texttt{linear\_combination} sobre \texttt{fp12 cyclotomic square} & \texttt{1622 -> 1886} \\
\texttt{linear\_combination} sobre \texttt{final exponentiation} & \texttt{587420 -> 678119} \\
\texttt{linear\_combination} sobre \texttt{pairing check} & \texttt{1682524 -> 1805233} \\
\bottomrule
\end{tabularx}

\section{Bloqueo actual y camino para destrabar el proyecto}

\subsection{Diagn\'ostico operativo actual}

El cuello de botella m\'as importante ya no es de plumbing, ni de sem\'antica circuital, ni de falta de
backend directo. La lane principal ya existe. El problema es \textbf{memoria} durante
\texttt{execute-wrapper-direct-prove-finalize}.

El s\'intoma documentado es:

\begin{quote}
\texttt{memory allocation of 268435456 bytes failed}
\end{quote}

Ese tama\~no coincide con un bloque de \(256\ \text{MiB}\) y es consistente con la hip\'otesis
de materializaci\'on pesada de polinomios/cosets en el extended domain para \texttt{k = 21}.

\subsection{Lo ya resuelto}

Las siguientes piezas ya fueron resueltas o materialmente mejoradas:

\begin{itemize}
  \item existe un \textbf{artefacto de setup m\'as rico} que evita rerun completo de \texttt{keygen\_pk(...)} en la wrapper lane;
  \item el repo mantiene un \textbf{patch local de \texttt{midnight-proofs}} con
  \texttt{BaseProvingKey}, \texttt{keygen\_pk\_base(...)},
  \texttt{create\_proof\_from\_base(...)},
  \texttt{create\_proof\_trace\_from\_base(...)} y
  \texttt{finalise\_proof\_from\_base\_trace(...)};
  \item el split \texttt{prove-trace} ya \textbf{funciona};
  \item el trace persistido ya migr\'o para guardar advice/instance polynomials en coefficient form;
  \item \texttt{HPolyKey} y \texttt{OpeningKey} ya redujeron duplicaci\'on de estado en varios puntos;
  \item existe chunking configurable para la parte de \texttt{h\_poly} vinculada a permutation constraints;
  \item el setup directo ya tiene al menos una corrida exitosa medida en \texttt{circom\_multiplier2}:
  \texttt{circuit\_k = 21}, \texttt{public\_input\_count = 1},
  \texttt{setup\_elapsed\_ms = 1554572} (aprox.\ 25m 54s).
\end{itemize}

\subsection{Lo parcialmente resuelto}

Hay mejoras que aislaron mejor el problema, pero no lo eliminaron:

\begin{itemize}
  \item el split \texttt{prove-trace}/\texttt{prove-finalize} reduce recomputaci\'on y permite diagn\'ostico m\'as fino;
  \item hay \textbf{fixed columns esparsas} y \textbf{lazy sigma cosets} por chunk en zonas relevantes;
  \item el tama\~no de chunk para la evaluaci\'on de \texttt{h\_poly} ya puede calibrarse desde CLI;
  \item el backend emite \textbf{checkpoints de finalize} m\'as finos, lo que permite localizar la \'ultima subfase exitosa antes del OOM.
\end{itemize}

\subsection{Lo no resuelto}

Lo que sigue abierto hoy es:

\begin{itemize}
  \item \texttt{execute-wrapper-direct-prove-finalize} todav\'ia aborta por memoria;
  \item el candidato principal sigue siendo el peso de estado alrededor de
  \texttt{compute\_h\_poly(...)} y la materializaci\'on de cosets extendidos;
  \item todav\'ia no est\'a decidido con evidencia suficiente si el siguiente recorte fuerte
  debe hacerse del lado \emph{pre-}\texttt{compute\_h\_poly} o \emph{post-}\texttt{compute\_h\_poly};
  \item por lo tanto, la lane directa existe pero no puede considerarse operacionalmente destrabada.
\end{itemize}

\section{Plan de acci\'on priorizado para destrabe}

El orden correcto, seg\'un el estado actual del repositorio y sus ADRs, es el siguiente.

\begin{enumerate}
  \item \textbf{Aislar el \'ultimo checkpoint exitoso de finalize.}\\
  Hay que usar la instrumentaci\'on actual para distinguir si el fallo ocurre antes o despu\'es de:
  \texttt{compute\_lagrange\_polys},
  sparse fixed cosets,
  permutation h key,
  y las subfases posteriores de opening/query.

  \item \textbf{Decidir con evidencia si el siguiente recorte debe ser pre-\texttt{compute\_h\_poly} o post-\texttt{compute\_h\_poly}.}\\
  El repositorio ya elimin\'o bastante duplicaci\'on wrapper-side. El siguiente paso no debe guiarse por intuici\'on:
  debe apoyarse en los logs de la instrumentaci\'on fina.

  \item \textbf{Seguir explotando sparse fixed columns, lazy sigma cosets y chunking.}\\
  La direcci\'on retenida actual es:
  materializar s\'olo lo necesario, lo m\'as tarde posible, y en porciones controlables.
  Esto debe profundizarse antes de intentar redise\~nos estructurales mayores.

  \item \textbf{Rerreproducir sistem\'aticamente con el split \texttt{prove-trace}/\texttt{prove-finalize}.}\\
  La experimentaci\'on debe regenerar artefactos cuando cambien formatos o supuestos de setup/trace/finalize.
  No conviene mezclar resultados obtenidos con artefactos obsoletos.

  \item \textbf{S\'olo despu\'es pasar al \texttt{h\_poly} speed follow-up.}\\
  El documento \texttt{docs/h-poly-followup-speed-plan.md} es expl\'icitamente posterior a la resoluci\'on del OOM.
  Antes de estabilizar memoria, optimizar velocidad es prematuro.
\end{enumerate}

En t\'erminos de prioridad ejecutiva, esto significa:

\begin{quote}
primero confiabilidad y estabilidad de memoria; despu\'es throughput.
\end{quote}

\section{Reproducibilidad}

\subsection{Fixtures de referencia}

El repositorio contiene dos fixtures principales para la lane actual:

\begin{itemize}
  \item \texttt{crates/wrapper-tests/fixtures/groth16/circom\_multiplier2/}
  como caso peque\~no y controlado;
  \item \texttt{crates/wrapper-tests/fixtures/groth16/semaphore/}
  como caso de integraci\'on m\'as application-shaped.
\end{itemize}

El fixture \texttt{circom\_multiplier2} corresponde a:

\begin{itemize}
  \item circuito \texttt{multiplier2.circom};
  \item entrada \(a = 3\), \(b = 11\);
  \item salida p\'ublica \(33\);
  \item artefactos \texttt{snarkjs} sobre la curva etiquetada como \texttt{bn128}, esto es, la familia BN254 usada por el repo.
\end{itemize}

\subsection{Comandos clave de CLI}

Los comandos m\'as relevantes para reproducir el estado actual son:

\begin{quote}
\texttt{cargo run -p wrapper-cli -- inspect-groth16-bundle ...}\\
\texttt{cargo run -p wrapper-cli -- plan-wrapper-job ...}\\
\texttt{cargo run -p wrapper-cli -- profile-layout}\\
\texttt{cargo run -p wrapper-cli -- execute-wrapper-direct-setup ...}\\
\texttt{cargo run -p wrapper-cli -- execute-wrapper-direct-prove ...}\\
\texttt{cargo run -p wrapper-cli -- execute-wrapper-direct-prove-trace ...}\\
\texttt{cargo run -p wrapper-cli -- execute-wrapper-direct-prove-finalize ...}\\
\texttt{cargo run -p wrapper-cli -- execute-wrapper-direct-verify ...}
\end{quote}

Para profiling de layout, la surface vigente incluye familias como:

\begin{itemize}
  \item \texttt{groth16};
  \item \texttt{pairing-terms};
  \item \texttt{public-inputs};
  \item \texttt{blocks};
  \item \texttt{outer}.
\end{itemize}

Para la experimentaci\'on del bloqueo actual, el comando m\'as sensible es
\texttt{execute-wrapper-direct-prove-finalize}, en combinaci\'on con el flag
\texttt{--h-poly-row-chunk-size} cuando se est\'a calibrando el tradeoff entre memoria y tiempo.

\subsection{Regla de higiene de artefactos}

La reproducci\'on fiel exige respetar la higiene de artefactos ya explicitada por el repositorio:

\begin{itemize}
  \item si cambia la producci\'on de setup, deben descartarse setup artifacts previos;
  \item si cambia la producci\'on o serializaci\'on de trace, deben descartarse traces previos;
  \item si cambia finalize o el formato del proof bundle, deben descartarse artifacts finales previos.
\end{itemize}

\section{Derivas documentales detectadas}

El relevamiento inicial detect\'o varias derivas documentales, y desde entonces
el repositorio fue corregido para que esas rutas vuelvan a ser navegables y la
documentaci\'on central apunte a referencias can\'onicas actuales.

\subsection{Archivo restaurado}

El repositorio ahora vuelve a incluir:

\begin{quote}
\texttt{docs/outer-prover-strategy-plan.md}
\end{quote}

Ese documento resume la estrategia vigente del prover outer y evita que el
estado de la lane directa tenga que reconstruirse s\'olo a partir de ADRs y
c\'odigo.

\subsection{Documento can\'onico de optimizaciones}

La referencia can\'onica actual para optimizaciones Midnight-locales es:

\begin{quote}
\texttt{docs/midnight-optimizations.md}
\end{quote}

El repositorio tambi\'en conserva una nota de compatibilidad en:

\begin{quote}
\texttt{docs/midnight-local-optimization-notes.md}
\end{quote}

de modo que enlaces hist\'oricos no queden rotos abruptamente.

\subsection{Compatibilidad para planes renombrados}

Las referencias a planes movidos bajo \texttt{docs/plans/} fueron corregidas en
la documentaci\'on central, y adem\'as el repositorio ahora expone archivos de
compatibilidad para rutas hist\'oricas como:

\begin{itemize}
  \item \texttt{docs/zk-email-integration-plan.md};
  \item \texttt{docs/plutus-aiken-integration-plan.md};
  \item \texttt{docs/cyclotomic-unitary-kernel-design.md}.
\end{itemize}

Eso permite que:

\begin{itemize}
  \item las referencias nuevas usen los paths can\'onicos bajo \texttt{docs/plans/};
  \item las referencias viejas sigan resolviendo sin ambig\"uedad.
\end{itemize}

\section{Conclusi\'on}

El proyecto ya alcanz\'o un punto importante: existe una lane directa real para envolver una slice
estrecha de verificaci\'on Groth16 BN254 dentro de un outer circuit Halo2/Midnight can\'onico.
La capa BN254, el pairing core y la reducci\'on Groth16 hoy est\'an implementados con suficiente
profundidad como para producir costos reales, optimizaciones medibles y cuellos de botella
aut\'enticos.

El principal bloqueo restante no es conceptual sino operativo: memoria en \texttt{prove-finalize}.
Por eso, el camino correcto para destrabar el proyecto no es ensanchar el alcance funcional,
sino completar el diagn\'ostico fino del estado alrededor de \texttt{compute\_h\_poly(...)} y seguir
reduciendo materializaci\'on simult\'anea de polinomios, cosets y estado de opening.

En s\'intesis: el proyecto ya sali\'o de la fase puramente experimental en arquitectura, pero todav\'ia
no sali\'o de la fase experimental en confiabilidad del prover directo.

\end{document}
